\section{Measuring the two noise scales from a single batch}
\label{sec:estimation}

Both $\tau_b$ and $\tau_w$ are recoverable from gradients a trainer already computes. Partition the
batch by prompt half and by rollout sub-group, and accumulate the four resulting microbatch
gradients separately:
\begin{equation}
  g_{ab} \;:=\; \text{mean gradient over prompt half } a \text{ using rollout sub-group } b,
  \qquad a, b \in \{0, 1\},
\end{equation}
each computed with its own sub-group's advantages. The two prompt halves are independent, and
given a prompt the two sub-groups are independent, so every term below is an inner product of
independent quantities:
\begin{align}
  \label{eq:estimators}
  \gbar^{\!\top} M \gbar &= \E\big[ g_{00}^{\!\top} M\, g_{11} \big],
  &&\text{different prompts, different responses,} \\
  \tr(M\Sb) &= \tfrac{P}{2}\Big( \E\big[ g_{a0}^{\!\top} M\, g_{a1} \big]
                                 - \E\big[ g_{00}^{\!\top} M\, g_{11} \big] \Big),
  &&\text{same prompts, different responses,} \nonumber \\
  \tr(M\Sw) &= \tfrac{P}{4}\, \E\big[ \| g_{a0} - g_{a1} \|_M^2 \big],
  &&\text{the prompt term cancels in the difference.} \nonumber
\end{align}

The cost is four gradient buffers rather than one, no extra responses, and no extra backward work
beyond accumulating into four accumulators instead of one -- which any trainer using gradient
accumulation over microbatches is already doing. Setting $M = I$ gives the practical form; setting
$M = F$ is obtained from measured KL rather than from materialising the Fisher.

\paragraph{Conditioning.} The obvious alternative, estimating the total variance and subtracting
the within-prompt part, is a difference of two large and highly correlated variances. On
homogeneous prompt pools it routinely returns a negative $\tau_b$ and an infinite $G^{*}$. The form
in \eqref{eq:estimators} differences two quantities of comparable size that share no sampling
noise, and it tracked the exact $\tau_b$ to within $4\%$ across pools spanning a factor of seven in
$G^{*}$ where the subtractive form failed outright.

\paragraph{Preconditioners.} An update is $\eta\, \Pi\, \ghat$ for a preconditioner $\Pi$, and Adam
is scale invariant, so both the drift model and the estimators must be read in the metric the
update actually moves in. In practice this means instrumenting the preconditioned gradient: all
results in \eqref{eq:estimators} carry over with $F \mapsto \Pi F \Pi$ under the standard
assumption that $\Pi$ varies slowly relative to batch-to-batch fluctuation.

\paragraph{Controlling drift.} Since $D \propto \eta^2$, a multiplicative controller
\begin{equation}
  \eta_{t+1} \;=\; \eta_t \cdot \mathrm{clip}\Big( \big( D^{*} / D_t \big)^{1/2},\ c^{-1},\ c \Big)
\end{equation}
corrects a mis-set step in a single update when the drift estimate is clean. Measured over runs on
transformers, this controller holds the median drift at its target with no detectable bias in log
space. Exponential smoothing of $D_t$, the obvious refinement, makes it worse: the lag biases the
realised drift low by roughly a factor of two.
